\documentclass[12pt,a4paper,twoside]{article}
\usepackage[utf-8]{inputenc}
\usepackage[romanian]{babel}
\usepackage{geometry}
\usepackage{hyperref}
\usepackage{graphicx}
\usepackage{listings}
\usepackage{xcolor}
\usepackage{fancyhdr}
\usepackage{array}
\usepackage{booktabs}
\usepackage{tikz}
\usepackage{float}
\usepackage{subcaption}
\usepackage{amsmath}

% Setare margini
\geometry{
    left=2.5cm,
    right=2.5cm,
    top=2.5cm,
    bottom=2.5cm,
    headheight=15pt,
    footskip=1cm
}

% Setare header/footer
\pagestyle{fancy}
\fancyhf{}
\fancyhead[L]{PCD T31 - SDD}
\fancyhead[R]{\thepage}
\fancyfoot[C]{Document de Proiectare --- Serviciu Web SOAP pentru Analiza Audio}

% Configurare listing-uri pentru cod
\lstset{
    language=C,
    basicstyle=\ttfamily\small,
    keywordstyle=\color{blue}\bfseries,
    commentstyle=\color{gray}\itshape,
    stringstyle=\color{red},
    breaklines=true,
    breakatwhitespace=true,
    tabsize=4,
    numbers=left,
    numberstyle=\tiny\color{gray},
    frame=single,
    backgroundcolor=\color{gray!5},
    captionpos=b
}

% Titlu și informații
\title{\textbf{PCD T31 --- Serviciu Web SOAP pentru Analiza Audio} \\
        \large{Software Design Document (SDD)}}
\author{Aplicație de Procesare Distribuită de Audio}
\date{Aprilie 2026}

\begin{document}

\maketitle

\begin{abstract}
    Acest document prezintă proiectarea detaliată a aplicației \textbf{PCD T31}, 
    un serviciu web SOAP pentru analiza și procesarea fișierelor audio/video. 
    Aplicația utilizează gSOAP pentru interfața de web service, LibrosaC pentru 
    procesarea DSP, libconfig pentru managementul configurației, și POSIX fork() 
    pentru paralelizare. Documentul acoperă arhitectura sistemului, componentele 
    principale, tipurile de date, protocoalele de comunicație, și detalii de implementare.
\end{abstract}

\newpage
\tableofcontents
\newpage

% =============================================================================
\section{Introducere}
% =============================================================================

\subsection{Scop și Context}

PCD T31 este o aplicație de procesare distribuită de date (Parallel Computing Distributed), 
fokusată pe \textbf{analiza spectrogramelor audio în timp real} printr-un serviciu web SOAP.

\textbf{Scopuri principale:}
\begin{itemize}
    \item Furnizarea unui serviciu web pentru analiza audio (spectrogram Mel)
    \item Suportul pentru procesare paralelă via fork()
    \item Configurație flexibilă la runtime prin libconfig
    \item Interacțiune client-server SOAP standardizată
    \item Suportul pentru fișiere audio și video
\end{itemize}

\subsection{Termeni și Definiții}

\begin{itemize}
    \item \textbf{SOAP} (Simple Object Access Protocol): Protocol pentru comunicație
          în servicii web.
    \item \textbf{gSOAP}: Toolkit open-source pentru generarea codul SOAP în C/C++.
    \item \textbf{LibrosaC}: Bibliotecă C pentru procesarea audio (spectrogram, MFCC, etc.).
    \item \textbf{Spectrograma}: Reprezentare vizuală a frecvențelor unui semnal în timp.
    \item \textbf{Mel scale}: Scală perceptuală a frecvențelor (utilizată în audio ML).
    \item \textbf{FFT}: Fast Fourier Transform, algoritm pentru calcul frecvență.
    \item \textbf{Hop length}: Pasul de deplasare în analiza timpului scurt (STFT).
    \item \textbf{Worker}: Proces copil creat via fork() pentru execuție izolată.
\end{itemize}

\subsection{Document Overview}

Structura acestui SDD urmează modelul standard IEEE 1016:

\begin{enumerate}
    \item \textbf{Sec. 2}: Descriere generală a sistemului
    \item \textbf{Sec. 3}: Structura arhitecturală
    \item \textbf{Sec. 4}: Componente și module software
    \item \textbf{Sec. 5}: Interfețe și protocoale
    \item \textbf{Sec. 6}: Fluxuri de date și control
    \item \textbf{Sec. 7}: Considerații de dependență și build
\end{enumerate}

% =============================================================================
\section{Descriere Generală a Sistemului}
% =============================================================================

\subsection{Perspectivă Sistem}

PCD T31 implementează o \textbf{arhitectură client-server} unde:

\begin{itemize}
    \item \textbf{Server (serverds):} Ascultă pe port 8080, acceptă cereri SOAP,
          spawneaza workers via fork() pentru procesare isolată
    \item \textbf{Clients (inetclient):} Se conectează la server, trimit cereri SOAP
          pentru analiza audio, așteaptă și consumă rezultate
    \item \textbf{Procesare:} Fiecare cerere este executată în procces copil pentru
          siguritate și paralelism
    \item \textbf{Configurație:} Runtime config din fișier \texttt{server.cfg} 
          (format libconfig)
\end{itemize}

\subsection{Restricții și Limitări}

\begin{table}[H]
    \centering
    \begin{tabular}{|l|l|}
        \toprule
        \textbf{Restricție} & \textbf{Valoare/Detaliu} \\
        \midrule
        Max workers simultani & 4 (configurable) \\
        Timeout SOAP & 30 secunde \\
        Max keep-alive per conexiune & 100 cereri \\
        Max path length & 512 bytes \\
        Backlog conexiuni TCP & 10 \\
        Standard C utilizat & C11 (\_POSIX\_C\_SOURCE=200809L) \\
        Port implicit & 8080 \\
        \bottomrule
    \end{tabular}
\end{table}

\subsection{Dependențe Externe}

\begin{itemize}
    \item \textbf{gsoap}: Generare cod SOAP C
    \item \textbf{libconfig}: Parsare fișiere configurație (libconfig/3.2 compatible)
    \item \textbf{librosacpp (LibrosaC)}: Bibliotecă DSP pentru spectrograme, MFCC, etc.
    \item \textbf{libpthread}: Suport POSIX threads (dacă folosit în viitor)
    \item \textbf{libm}: Bibliotecă matematică (pentru pow, log, etc.)
    \item \textbf{POSIX syscalls}: write(), fork(), waitpid(), open(), read(), close(), etc.
\end{itemize}

% =============================================================================
\section{Arhitectura Sistemului}
% =============================================================================

\subsection{Diagrama Conceptuală}

\begin{verbatim}
┌─────────────────────────────────────────────────────────────────┐
│                    CLIENT APPLICATION (inetclient)            │
├─────────────────────────────────────────────────────────────────┤
│  ┌─────────────────────────────────────────────────────────┐   │
│  │ SOAP Proxy (generated by gSOAP)                         │   │
│  │  - soap_call_ns__AnalyzeAudio()                         │   │
│  │  - HTTP over TCP                                        │   │
│  └─────────────────────────────────────────────────────────┘   │
└────────────────────┬──────────────────────────────────────────┬─┘
                     │ HTTP POST (SOAP envelope)                │ 
                     │ Port 8080                                │
┌────────────────────┴──────────────────────────────────────────┴─┐
│                   SERVER APPLICATION (serverds)               │
├──────────────────────────────────────────────────────────────────┤
│ ┌────────────────────────────────────────────────────────────┐  │
│ │ SOAP Service Handler (main process)                      │  │
│ │  - soap_serve() loop                                      │  │
│ │  - Accept conexiuni TCP                                   │  │
│ │  - Parse SOAP envelope                                    │  │
│ │  - fork() new worker process                              │  │
│ └────────────────────────────────────────────────────────────┘  │
│                           │                                     │
│    ┌──────────────────────┼──────────────────────┐             │
│    │                      │                      │             │
│ ┌──┴────┐  ┌──────┐  ┌────┴─┐  ┌──────┐        │             │
│ │Worker1│  │Worker2│  │Worker3│  │Worker4│        │             │
│ │(fork) │  │(fork) │  │(fork) │  │(fork) │        │             │
│ └───┬───┘  └──┬───┘  └───┬───┘  └──┬───┘        │             │
│     │         │          │         │            │             │
│     └─────────┼──────────┼─────────┘            │             │
│               │          │                      │             │
│     ┌─────────▼──────────▼──────────┐           │             │
│     │ Config Loader (libconfig)     │           │             │
│     │  - port, max_workers, etc.    │           │             │
│     └─────────────────────────────────┘           │             │
│                                                   │             │
│     ┌──────────────────────────────────┐         │             │
│     │ Processing Module (LibrosaC)     │         │             │
│     │  - Mel spectrogram               │         │             │
│     │  - File I/O (output/)            │         │             │
│     └──────────────────────────────────┘         │             │
└──────────────────────────────────────────────────┴──────────────┘
\end{verbatim}

\subsection{Stratificarea Arhitecturii}

Aplicația este organizată în \textbf{4 straturi}:

\begin{enumerate}
    \item \textbf{Transport Layer}: SOAP/HTTP, geolocalizări gSOAP
    \item \textbf{Service Layer}: SOAP handlers, request dispatch
    \item \textbf{Processing Layer}: DSP engine (LibrosaC), audio I/O
    \item \textbf{Configuration Layer}: Runtime config (libconfig)
\end{enumerate}

\subsection{Modelul Concurenței}

\textbf{Forking-based parallelism:}

\begin{itemize}
    \item Procesul principal (serverds) ascultă pe portul SOAP
    \item La fiecare cerere SOAP, se crează un \textbf{worker process} via \texttt{fork()}
    \item Workerul procesează solicitarea în spațiu de memorie izolat
    \item Workerul se termină; procesul principal așteptă via \texttt{waitpid()}
    \item Max concurrent workers = \texttt{max\_workers} din config
\end{itemize}

\textbf{Avantaje:}
\begin{itemize}
    \item Izolare robustă (crash worker $\neq$ crash server)
    \item Simplitate (nu trebuie mutex, conditions, etc.)
    \item Portabilitate (POSIX standard)
\end{itemize}

% =============================================================================
\section{Componente și Module Software}
% =============================================================================

\subsection{Modulul Server (serverds)}

\subsubsection{Fișiere Componente}
\begin{itemize}
    \item \texttt{src/server.c} --- Logica server principal
    \item \texttt{include/server.h} --- Tipuri și constante server
    \item \texttt{src/soapds.c} --- Implementare SOAP handlers
\end{itemize}

\subsubsection{Responsabilități}

\begin{itemize}
    \item Inițializare SOAP service (gSOAP runtime)
    \item Parsare argumentelor CLI (-c, -p, -v, -h)
    \item Încărcare configurație din fișier
    \item Setup signal handlers (SIGTERM pentru graceful shutdown)
    \item Accept conexiuni TCP
    \item Dispatch cereri SOAP la handlers
    \item Fork() workers pentru procesare isolată
    \item Cleanup și shutdown
\end{itemize}

\subsubsection{Interfața Publică}

\begin{lstlisting}
/* Tipul configurației server-ului */
typedef struct {
    int  port;                    /* Port SOAP (ex: 8080) */
    int  max_workers;             /* Max procese worker (ex: 4) */
    int  log_level;               /* 0=silent, 1=info, 2=debug */
    char output_dir[MAX_PATH_LEN]; /* Directorul output */
    char soap_endpoint[MAX_PATH_LEN]; /* URL endpoint */
} server_config_t;

/* Constante */
#define DEFAULT_PORT         8080
#define DEFAULT_CONFIG_PATH  "config/server.cfg"
#define MAX_PATH_LEN         512
#define MAX_WORKERS_DEFAULT  4
#define SOAP_BACKLOG         10
#define SOAP_TIMEOUT_S       30
#define SOAP_MAX_KEEP_ALIVE  100
\end{lstlisting}

\subsection{Modulul SOAP Handlers (soapds.c)}

\subsubsection{Responsabilități}

\begin{itemize}
    \item Implementarea operațiunilor SOAP (C binding-uri)
    \item Mapare cereri HTTP SOAP $\rightarrow$ C structs
    \item Invocarea logicii de procesare
    \item Construire răspunsuri SOAP
    \item Gestiunea erorilor
\end{itemize}

\subsubsection{Operații SOAP Expuse}

\begin{enumerate}
    \item \textbf{AnalyzeAudio}: Calculează spectrograma Mel pentru un fișier audio
    \item \textbf{GetJobStatus}: (Milestone viitor) Query status job async
    \item \textbf{ExtractFrame}: (Milestone viitor) Extrage frame din video
\end{enumerate}

\subsection{Modulul Processing (processing.c / processing.h)}

\subsubsection{Fișiere Componente}
\begin{itemize}
    \item \texttt{src/processing.c} --- Implementare DSP
    \item \texttt{include/processing.h} --- Interfață publică
\end{itemize}

\subsubsection{Responsabilități}

\begin{itemize}
    \item Inițializare LibrosaC (dacă necesar)
    \item Citire fișier audio din FS
    \item Aplicare resample (target\_sr)
    \item Calculare STFT (FFT, window, overlap)
    \item Conversie la Mel scale (n\_fft, hop\_length, n\_mels)
    \item Normalizare și log-compresie
    \item Scriere spectrogram binar în output/
    \item Returnare metadate (sample rate, frames, durata)
\end{itemize}

\subsubsection{Tipuri de Date}

\begin{lstlisting}
/* Cerere de procesare */
typedef struct {
    char input_path[MAX_PATH_LEN];   /* Cale intrare audio/video */
    char output_path[MAX_PATH_LEN];  /* Cale ieșire spectrogram */
    int  target_sr;                  /* Target sample rate (Hz) */
    int  n_fft;                      /* FFT size */
    int  hop_length;                 /* Hop length (frames) */
    int  n_mels;                     /* Meluri benzi */
} proc_request_t;

/* Rezultat procesare */
typedef struct {
    int    status;                       /* 0=OK, -1=error */
    char   error_msg[256];               /* Mesaj eroare */
    char   output_path[MAX_PATH_LEN];    /* Cale output */
    int    sample_rate;                  /* Sample rate actual */
    int    n_samples;                    /* Număr samples */
    double duration_s;                   /* Durata (secunde) */
    int    n_mels;                       /* Mel benzi output */
    int    n_frames;                     /* Cadre timp output */
} proc_result_t;

int process_spectrogram(const proc_request_t*req, 
                        proc_result_t*res);
\end{lstlisting}

\textbf{Valori Implicite:}

\begin{table}[H]
    \centering
    \begin{tabular}{|l|r|}
        \toprule
        \textbf{Parametru} & \textbf{Valoare Implicită} \\
        \midrule
        Sample Rate & 22050 Hz \\
        N FFT & 2048 \\
        Hop Length & 512 \\
        N Mels & 128 \\
        \bottomrule
    \end{tabular}
\end{table}

\subsection{Modulul Config Loader (config\_loader.c / .h)}

\subsubsection{Responsabilități}

\begin{itemize}
    \item Parse fișier \texttt{.cfg} (format libconfig)
    \item Extract parametri server (port, max\_workers, log\_level, etc.)
    \item Validare valori (port în range 1024-65535, etc.)
    \item Alocare memorie pentru strings (output\_dir, endpoint)
    \item Cleanup resurse (config\_free)
\end{itemize}

\subsubsection{Interfață Publică}

\begin{lstlisting}
int  config_load(const char*path, server_config_t*out);
void config_free(server_config_t*cfg);
\end{lstlisting}

\subsubsection{Format Fișier Config}

\begin{lstlisting}
# PCD T31 - Server Configuration File
# Format: libconfig

server:
{
    port = 8080;
    max_workers = 4;
    output_dir = "./output";
    soap_endpoint = "http://localhost:8080/soap";
    log_level = 1;
};

processing:
{
    default_sample_rate = 22050;
    default_n_fft       = 2048;
    default_hop_length  = 512;
    default_n_mels      = 128;
};
\end{lstlisting}

\subsection{Modulul Client (inetclient)}

\subsubsection{Fișiere Componente}
\begin{itemize}
    \item \texttt{src/client.c} --- Logica client SOAP
\end{itemize}

\subsubsection{Responsabilități}

\begin{itemize}
    \item Parse argumente CLI (-s server, -f file, -m mels, -r sr, -h help)
    \item Citire variabilă mediu \texttt{PCD\_SERVER} (overrides default)
    \item Construire cerere SOAP AnalyzeAudio
    \item Apel soap\_call\_ns\_\_AnalyzeAudio()
    \item Afișare rezultate (status, output path, metadate)
    \item Cleanup SOAP context
\end{itemize}

\subsubsection{Utilizare}

\begin{lstlisting}
# Apel cu argumente
./inetclient -s http://localhost:8080/soap \
             -f audio.wav -m 128 -r 22050

# Cu variabilă mediu
PCD_SERVER=http://server2.com:9090/soap \
  ./inetclient -f audio.wav
\end{lstlisting}

\subsection{Module Demo (demo\_milestone1, inetds2, inetsample2)}

Aplicații de demonstrație și test pentru validare a componentelor.

% =============================================================================
\section{Interfețe și Protocoale}
% =============================================================================

\subsection{Protocolul SOAP}

\subsubsection{Generare și Compilare}

Fișierul \texttt{src/proto.h} conține definițiile interface-ului SOAP:

\begin{lstlisting}
/**
 * gSOAP namespace map
 * Serviciu: ServiciulVideoSpectrograma
 * Endpoint: http://localhost:8080/soap
 * Schema namespace: urn:ServiciulVideoSpectrograma
 */

//gSOAP ns serviciu nume: ServiciulVideoSpectrograma
//gSOAP ns serviciu style: document
//gSOAP ns serviciu encoding: literal
//gSOAP ns serviciu location: http://localhost:8080/soap
//gSOAP ns schema namespace: urn:ServiciulVideoSpectrograma
\end{lstlisting}

Generare cod:
\begin{lstlisting}
soapcpp2 -c -S -Iinclude -Isrc src/proto.h -d src/
\end{lstlisting}

Aceasta produce:
\begin{itemize}
    \item \texttt{src/soapC.c} --- Marshaling/unmarshaling funcții
    \item \texttt{src/soapServer.c} --- Server dispatch
    \item \texttt{src/soapClient.c} --- Client proxy (nu se folosește direct, doar linkare)
    \item \texttt{src/soapH.h} --- Header generat cu types
    \item \texttt{src/soap*.h} --- Helper headers
\end{itemize}

\subsubsection{Operația AnalyzeAudio}

\textbf{Cerere HTTP SOAP:}

\begin{lstlisting}
POST /soap HTTP/1.1
Host: localhost:8080
Content-Type: text/xml; charset=utf-8
Content-Length: XXX
SOAPAction: ""

<?xml version="1.0" encoding="UTF-8"?>
<SOAP-ENV:Envelope xmlns:SOAP-ENV="..." 
                   xmlns:ns="urn:ServiciulVideoSpectrograma">
  <SOAP-ENV:Body>
    <ns:AnalyzeAudio>
      <inputFilePath>/path/to/audio.wav</inputFilePath>
      <targetSampleRate>22050</targetSampleRate>
      <nMels>128</nMels>
      <nFft>2048</nFft>
      <hopLength>512</hopLength>
    </ns:AnalyzeAudio>
  </SOAP-ENV:Body>
</SOAP-ENV:Envelope>
\end{lstlisting}

\textbf{Răspuns HTTP SOAP:}

\begin{lstlisting}
HTTP/1.1 200 OK
Content-Type: text/xml; charset=utf-8
Content-Length: YYY

<?xml version="1.0" encoding="UTF-8"?>
<SOAP-ENV:Envelope xmlns:SOAP-ENV="..." 
                   xmlns:ns="urn:ServiciulVideoSpectrograma">
  <SOAP-ENV:Body>
    <ns:AnalyzeAudioResponse>
      <statusCode>0</statusCode>
      <statusMessage>OK</statusMessage>
      <resultFilePath>./output/spec_12345.bin</resultFilePath>
      <sampleRate>22050</sampleRate>
      <durationSeconds>45.5</durationSeconds>
      <nMels>128</nMels>
      <nFrames>1234</nFrames>
    </ns:AnalyzeAudioResponse>
  </SOAP-ENV:Body>
</SOAP-ENV:Envelope>
\end{lstlisting}

\subsection{Configurația Libconfig}

Format fișier: grup server cu sub-parametri.

\begin{itemize}
    \item \textbf{port}: Port TCP pentru SOAP (int, 1024-65535)
    \item \textbf{max\_workers}: Max procese worker (int, $\geq 1$)
    \item \textbf{output\_dir}: Directory ieșire spectrogram (string)
    \item \textbf{soap\_endpoint}: URL endpoint (string, informativ)
    \item \textbf{log\_level}: 0=silent, 1=info, 2=debug
\end{itemize}

\subsection{Interfața File I/O}

\textbf{Intrări:}
\begin{itemize}
    \item Fișiere audio: WAV, MP3, FLAC, OGG (suportate de LibrosaC)
    \item Fișiere video: MP4, AVI, MKV (audio extract)
\end{itemize}

\textbf{Ieșiri:}
\begin{itemize}
    \item Fișiere spectrogram: Format binar (definit de LibrosaC)
    \item Locație: \texttt{./output/} (configurable)
    \item Nomenclare: \texttt{spec\_<timestamp>.bin} sau similară
\end{itemize}

% =============================================================================
\section{Fluxuri de Date și Control}
% =============================================================================

\subsection{Fluxul Request/Response}

\begin{verbatim}
┌─────────────────┐
│   Client Ready  │
└────────┬────────┘
         │
         │ soap_call_ns__AnalyzeAudio(&req, &res)
         │
         ▼
┌─────────────────────────┐
│  Trimite SOAP Request   │
│   HTTP POST             │
└────────┬────────────────┘
         │
         ▼
┌──────────────────────────┐
│  Server Receives         │
│  HTTP POST               │
└────────┬─────────────────┘
         │
         ▼
┌──────────────────────────┐
│  gSOAP Parse SOAP Body   │
│  Unmarshal to C struct   │
└────────┬─────────────────┘
         │
         ▼
┌──────────────────────────┐
│  Call Handler Function   │
│  (ns__AnalyzeAudio)      │
└────────┬─────────────────┘
         │
         ▼
┌──────────────────────────┐
│  fork() Worker Process   │
└────────┬─────────────────┘
         │
    ┌────┴─────────────────────────┐
    │                              │
    ▼ (Parent)                     ▼ (Child)
┌──────────────┐          ┌──────────────────────┐
│ waitpid()    │          │ process_spectrogram()│
│              │          │ (LibrosaC)           │
└──────────────┘          └──────┬───────────────┘
                                 │
                                 ▼
                          ┌──────────────────┐
                          │ Write output file│
                          │ to ./output/     │
                          └──────┬───────────┘
                                 │
                                 ▼
                          ┌──────────────────┐
                          │ exit(0)          │
                          └──────────────────┘
         │
         ▼
┌──────────────────────────┐
│  Marshal Result to SOAP  │
│  Build Response Envelope │
└────────┬─────────────────┘
         │
         ▼
┌──────────────────────────┐
│  HTTP 200 OK             │
│  Send SOAP Response      │
└────────┬─────────────────┘
         │
         ▼
┌──────────────────────────┐
│  Client Receives         │
│  Unmarshals Response     │
└────────┬─────────────────┘
         │
         ▼
┌──────────────────────────┐
│  Client Displays Result  │
└──────────────────────────┘
\end{verbatim}

\subsection{Fluxul Intern Processing}

\begin{enumerate}
    \item \textbf{Parsare cerere:} Extract input\_path, target\_sr, n\_mels, etc.
    \item \textbf{Validare fișier:} Check dacă input\_path există, readable
    \item \textbf{Inițializare LibrosaC:} Setare rate, FFT size, etc.
    \item \textbf{Citire audio:} Load fișier, decode (WAV/MP3/etc)
    \item \textbf{Resample:} Dacă target\_sr $\neq$ native rate
    \item \textbf{STFT:} Aplicare window (Hann), FFT, compute magnitude
    \item \textbf{Mel Filter Bank:} Mapare frecvență $\rightarrow$ Mel scale
    \item \textbf{Log Compression:} Log10(1 + mel\_spec)
    \item \textbf{Scriere output:} Binar format (float32 matrix)
    \item \textbf{Calcul metadate:} n\_samples, duration, n\_frames
    \item \textbf{Return result:} Status=0, metadate populate
\end{enumerate}

\subsection{State Diagram --- Server Main Process}

\begin{verbatim}
┌──────────┐
│  START   │
└────┬─────┘
     │
     ▼
┌────────────────────┐
│ Parse CLI args     │
│ Load config file   │
├────────────────────┤
│  config_load()     │
└────┬───────────────┘
     │
     ▼
┌────────────────────┐
│ Setup signal       │
│ handlers (SIGTERM) │
└────┬───────────────┘
     │
     ▼
┌────────────────────────────┐
│ gSOAP Init                 │
│ Listen on TCP port         │
└────┬───────────────────────┘
     │
     ▼
┌────────────────────────────┐
│  g_running == 1 ?          │
├────────────────────────────┤
│ /          \               │
│ NO         YES (continue)  │
│ │          │               │
└─┼──────────┼───────────────┘
  │          │
  │          ▼
  │     ┌──────────────┐
  │     │ soap_serve() │
  │     │ (blocking)   │
  │     │ Accept conn  │
  │     │ Handle SOAP  │
  │     └──────┬───────┘
  │            │
  │            ▼
  │     ┌──────────────────────┐
  │     │ fork() worker        │
  │     │ or direct handle     │
  │     └──────┬───────────────┘
  │            │
  │            ▼
  │     ┌──────────────────────┐
  │     │ Process request      │
  │     │ (worker or thread)   │
  │     └──────┬───────────────┘
  │            │
  │            ▼
  │     ┌──────────────────────┐
  │     │ Loop back            │
  │     │ (next conn or SIGTERM)
  │     └──────┬───────────────┘
  │            │
  │            └──────► Loop to g_running check
  │
  ▼
┌────────────────────────┐
│ Cleanup resources      │
│ soap_destroy()         │
│ config_free()          │
└────┬───────────────────┘
     │
     ▼
┌────────────────────────┐
│  EXIT (0)              │
└────────────────────────┘
\end{verbatim}

% =============================================================================
\section{Considerații Tehnologice și Dependență Build}
% =============================================================================

\subsection{Build System --- Makefile}

\begin{lstlisting}
CC      = gcc
CFLAGS  = -Wall -Wextra -Wpedantic -Wformat=2 -Wshadow \
          -Wstrict-prototypes -Wmissing-prototypes       \
          -Wno-unused-parameter                          \
          -std=c11 -D_POSIX_C_SOURCE=200809L             \
          -Iinclude -Isrc

LIBS_SOAP    = -lgsoap
LIBS_CONFIG  = -lconfig
LIBS_LIBROSA = -lrosacpp
LIBS_THREAD  = -lpthread
LIBS_MATH    = -lm
\end{lstlisting}

\textbf{Ținte build:}

\begin{itemize}
    \item \texttt{make all}: Compilare toate executabile
    \item \texttt{make serverds}: Compilare SOAP server
    \item \texttt{make inetclient}: Compilare SOAP client
    \item \texttt{make demo\_milestone1}: Compilare demo app
    \item \texttt{make soap\_gen}: Regenerare cod SOAP din proto.h
    \item \texttt{make tidy}: Run clang-tidy (static analysis)
    \item \texttt{make clean}: Ștergere binare și fișiere generate
\end{lstlisting}

\subsection{Dependențe Externe --- Package Management}

\begin{table}[H]
    \centering
    \begin{tabular}{|l|l|l|}
        \toprule
        \textbf{Bibliotecă} & \textbf{Versiune} & \textbf{Scop} \\
        \midrule
        gsoap & 2.8+ & Generare SOAP C code \\
        libconfig & 1.7+ & Parse \texttt{.cfg} files \\
        librosacpp & 0.9+ & DSP (Mel spectrograma) \\
        libpthread & POSIX & Threading (viitor) \\
        libm & Standard & Math functions \\
        libc (glibc) & 2.31+ & POSIX syscalls \\
        \bottomrule
    \end{tabular}
\end{table}

\subsection{Flag-uri Compiler și Opțiuni}

\textbf{Optimizare și Debugging:}

\begin{lstlisting}
# Release build
make DEBUG=0    # -O2 -DNDEBUG

# Debug build
make DEBUG=1    # -g -O0 -DDEBUG
\end{lstlisting}

\textbf{Flags de Avertisment:}

\begin{itemize}
    \item \texttt{-Wall -Wextra -Wpedantic}: Toate avertismentele standard
    \item \texttt{-Wformat=2}: Verificare printf/scanf format strings
    \item \texttt{-Wshadow}: Variabile shadowing
    \item \texttt{-Wstrict-prototypes}: Functii fără prototype
\end{itemize}

\textbf{Standard C:}

\begin{itemize}
    \item \texttt{-std=c11}: C11 standard (variadic macros, generics, etc.)
    \item \texttt{-D\_POSIX\_C\_SOURCE=200809L}: POSIX 2008 symbols
\end{itemize}

\subsection{Generare SOAP Bindings}

\begin{lstlisting}
# Generate C stubs from proto.h
soapcpp2 -c -S -Iinclude -Isrc src/proto.h -d src/

# Opțiuni:
# -c       : Generate C code (not C++)
# -S       : Generate server code (dispatch)
# -Iinclude : Include path untuk gSOAP headers
# -d src/  : Output directory
\end{lstlisting}

\textbf{Fișiere generate:}

\begin{itemize}
    \item \texttt{src/soapC.c}: Core marshaling funcții
    \item \texttt{src/soapServer.c}: Service implementation skeletons
    \item \texttt{src/soapClient.c}: Client proxy (pentru reference)
    \item \texttt{src/soapH.h}: Typedef-uri și prototypes
    \item \texttt{src/soapStub.h}: Struct definitions
\end{itemize}

\subsection{Analiza Statică cu clang-tidy}

\begin{lstlisting}
make tidy    # Run clang-tidy sur file-urile critice
\end{lstlisting}

Verifică:
\begin{itemize}
    \item Memory leaks (valgrind integration)
    \item Use-after-free
    \item Buffer overflows
    \item Dead code
    \item Style violations
\end{itemize}

% =============================================================================
\section{Managementul Erorilor și Logging}
% =============================================================================

\subsection{Niveluri de Logging}

\begin{table}[H]
    \centering
    \begin{tabular}{|c|l|l|}
        \toprule
        \textbf{Level} & \textbf{Nume} & \textbf{Descriere} \\
        \midrule
        0 & SILENT & Nici un output \\
        1 & INFO & Mesaje informative (default) \\
        2 & DEBUG & Mesaje verbose pentru debugging \\
        \bottomrule
    \end{tabular}
\end{table}

\textbf{Configurare:}

\begin{lstlisting}
# În server.cfg
server:
{
    log_level = 1;  # 0=silent, 1=info, 2=debug
};
\end{lstlisting}

\subsection{Coduri Eroare}

\begin{table}[H]
    \centering
    \begin{tabular}{|c|l|}
        \toprule
        \textbf{Cod} & \textbf{Semnificație} \\
        \midrule
        0 & Succes \\
        -1 & Eroare generică (vezi error\_msg) \\
        -2 & Fișier nu găsit \\
        -3 & Permisiuni insuficiente \\
        -4 & Out of memory \\
        -5 & Invalid parametru \\
        \bottomrule
    \end{tabular}
\end{table}

\subsection{Strategie Cleanup}

\textbf{Resurse în SOAP request/response:}

\begin{itemize}
    \item \texttt{soap\_end()} --- Cleanup după fiecare cerere
    \item \texttt{soap\_destroy()} --- Cleanup la shutdown
    \item Pointeri alare: \texttt{soap\_malloc()}, \texttt{soap\_strdup()}
\end{itemize}

\textbf{Resurse în processing:}

\begin{itemize}
    \item Fișiere: \texttt{open()}, \texttt{read()}, \texttt{close()}
    \item Bufere: Alocare cu \texttt{malloc()}, free cu \texttt{free()}
\end{itemize}

% =============================================================================
\section{Seguritate și Robustețe}
% =============================================================================

\subsection{Input Validation}

\begin{itemize}
    \item \textbf{Path traversal:} Validează paths (nu permite \texttt{../} escape)
    \item \textbf{Buffer overflow:} Folosire \texttt{strncpy()}, size limits
    \item \textbf{Integer overflow:} Verificare ranges pentru n\_fft, n\_mels, etc.
    \item \textbf{SOAP injection:} gSOAP sanitizează XML automatically
\end{itemize}

\subsection{Process Isolation via fork()}

\begin{itemize}
    \item Fiecare worker ruleaza în spațiu de memorie separat
    \item Crash worker $\neq$ crash server
    \item Resurse cleanup automatice pe exit process
    \item Semnale (SIGSEGV, SIGABRT) nu afectează parent
\end{itemize}

\subsection{Signal Handling}

\begin{itemize}
    \item \textbf{SIGTERM}: Graceful shutdown, cleanup
    \item \textbf{SIGCHLD}: (Opțional) Reap zombie workers
    \item \textbf{SIGPIPE}: Ignore (previne crash pe broken socket)
    \item Signal handlers sunt async-signal-safe (doar \texttt{write()}, etc.)
\end{itemize}

\subsection{Limitări Resurse}

\begin{itemize}
    \item \textbf{max\_workers:} Limit concurrent processes
    \item \textbf{SOAP\_TIMEOUT\_S:} Timeout 30s pe operație
    \item \textbf{MAX\_KEEP\_ALIVE:} Limit reuse conexiuni
    \item \textbf{MAX\_PATH\_LEN:} 512 bytes limit paths
\end{itemize}

% =============================================================================
\section{Testare și Validare}
% =============================================================================

\subsection{Strategia de Testare}

\begin{enumerate}
    \item \textbf{Unit Tests}: Test componente individual (processing, config\_loader)
    \item \textbf{Integration Tests}: Test SOAP client-server interaction
    \item \textbf{Load Tests}: Test concurrent workers, stress testing
    \item \textbf{Security Tests}: Path traversal, buffer overflow, fuzzing
    \item \textbf{Static Analysis}: clang-tidy, code review
\end{enumerate}

\subsection{Scripturi Test}

\begin{itemize}
    \item \texttt{test\_milestone1.sh}: Executa scenarii test
    \item \texttt{run\_demo.sh}: Run demo application
    \item \texttt{setup.sh}, \texttt{setup\_complet.sh}: Environment setup
\end{itemize}

\subsection{Cazuri de Test Relevante}

\begin{enumerate}
    \item \textbf{Happy Path}: Client-server request/response normal
    \item \textbf{Missing File}: Input file nu există
    \item \textbf{Invalid Params}: Parameters invalid (n\_mels=0, sr=0, etc.)
    \item \textbf{Concurrent Requests}: Multiple clients simultaneous
    \item \textbf{Server Shutdown}: SIGTERM graceful shutdown
    \item \textbf{Worker Crash}: Verify server survives worker crash
    \item \textbf{Disk Full}: Output directory full, error handling
\end{enumerate}

% =============================================================================
\section{Performance și Scalabilitate}
% =============================================================================

\subsection{Metrici de Performance}

\begin{itemize}
    \item \textbf{Latență request:} Timp de la client POST până la response primire
    \item \textbf{Throughput:} Cereri/sec procesate concurrent
    \item \textbf{CPU Usage}: Per worker, total server
    \item \textbf{Memory Usage}: Peak RSS per worker
    \item \textbf{I/O Efficiency}: Disk read/write speeds
\end{itemize}

\subsection{Bottleneck-uri Identificate}

\begin{enumerate}
    \item \textbf{LibrosaC DSP}: Calculul STFT și Mel filtri sunt CPU-intensive
    \item \textbf{Disk I/O}: Write output spectrogram la disk
    \item \textbf{Network Latency}: SOAP over HTTP
    \item \textbf{Process forking}: Overhead syscall fork()
\end{enumerate}

\subsection{Optimizări Posibile (Viitor)}

\begin{itemize}
    \item \textbf{Thread Pool}: Înlocuire fork() cu thread pool (pthreads)
    \item \textbf{Async I/O}: Folosire aio\_*, epoll() pentru I/O async
    \item \textbf{Caching}: Cache spectrograme deja calculate
    \item \textbf{Compression}: Comprese output spectrogram (gzip)
    \item \textbf{REST API}: REST + JSON alternativ la SOAP (lightweight)
\end{itemize}

% =============================================================================
\section{Desfășurare și Operaționale}
% =============================================================================

\subsection{Cerinți Sistem}

\begin{itemize}
    \item \textbf{OS}: Linux (Kernel 3.10+), macOS 10.12+
    \item \textbf{Compiler}: GCC 5.0+, Clang 6.0+
    \item \textbf{RAM}: Min 512MB, recomandat 2GB+
    \item \textbf{Disk}: Min 100MB (pentru build + output)
    \item \textbf{Network}: TCP/IP stack (pentru SOAP)
\end{itemize}

\subsection{Instalare Dependencies}

\subsubsection{Debian/Ubuntu}
\begin{lstlisting}
sudo apt-get install gcc make git
sudo apt-get install gsoap libconfig-dev
sudo apt-get install libpthread-stubs0-dev libm librosacpp-dev
\end{lstlisting}

\subsubsection{RedHat/CentOS}
\begin{lstlisting}
sudo yum install gcc make git
sudo yum install gsoap libconfig-devel
sudo yum install pthread-devel libm librosacpp-devel
\end{lstlisting}

\subsection{Build și Desfășurare}

\begin{lstlisting}
# 1. Clone sau cd workspace
cd pcd-t31

# 2. Generate SOAP bindings
make soap_gen

# 3. Build all
make all

# 4. Verify executabile
ls -la serverds inetclient demo_milestone1

# 5. Start server
./serverds -c config/server.cfg -p 8080

# 6. (În alt terminal) Run client
./inetclient -f audio.wav -m 128

# 7. Check output
ls -la output/
\end{lstlisting}

\subsection{Monitoring și Diagnostică}

\begin{itemize}
    \item \textbf{ps aux}: Verifică running processes (serverds, workers)
    \item \textbf{netstat -tuln}: Verifică port 8080 listening
    \item \textbf{top}: Monitor CPU/memory per process
    \item \textbf{strace}: Trace syscalls (debugging)
    \item \textbf{valgrind}: Memory profiling, leak detection
\end{itemize}

\subsection{Graceful Shutdown}

\begin{lstlisting}
# Trimite SIGTERM la serverds
kill -TERM <PID>

# Server va:
# 1. Set g_running = 0
# 2. Close accept() loop
# 3. Wait for active workers (waitpid)
# 4. Cleanup resources
# 5. Exit cleanly
\end{lstlisting}

% =============================================================================
\section{Referințe și Resurse}
% =============================================================================

\begin{itemize}
    \item \textbf{gSOAP}: \url{https://www.genivia.com/products.html}
    \item \textbf{libconfig}: \url{https://hyperrealm.github.io/libconfig/}
    \item \textbf{LibrosaC}: Audio DSP library (C binding)
    \item \textbf{POSIX Standard}: \url{https://pubs.opengroup.org/onlinepubs/}
    \item \textbf{IEEE 1016-2009}: Software Design Document standard
    \item \textbf{HTTP/1.1 RFC}: \url{https://tools.ietf.org/html/rfc7230}
    \item \textbf{SOAP 1.2 W3C}: \url{https://www.w3.org/TR/soap12/}
\end{itemize}

% =============================================================================
\section{Concluzie}
% =============================================================================

PCD T31 este o aplicație de procesare distribuită bine structurată, implementând:

\begin{itemize}
    \item Serviciu web SOAP standard pentru analiza audio
    \item Arhitectură forking-based pentru paralelism și robustețe
    \item Configurație flexibilă via libconfig
    \item Componente modulare cu responsabilități clare
    \item Security și error handling corespunzători
    \item Build system simplu și reproducibil
\end{itemize}

\textbf{Puncte forte:}
\begin{itemize}
    \item Izolare proceselor via fork()
    \item SOAP standard, generare automată cod gSOAP
    \item Logging și error handling robusti
    \item Suport multi-format audio/video (LibrosaC)
\end{itemize}

\textbf{Arii de îmbunătățire:}
\begin{itemize}
    \item Migrare la thread pool (eficiență memorie)
    \item Async I/O pentru scalabilitate mai bună
    \item REST API alternativ (lightweight clients)
    \item Caching spectrograme calculate anterior
\end{itemize}

Aplicația este gata pentru deployment în mediu de producție cu monitoring și scaling corespunzător.

\appendix

% =============================================================================
\section{Diagrama Dependințelor Modului}
% =============================================================================

\begin{verbatim}
┌──────────────────────────────────────────────────────────────────┐
│                     EXECUTABILE                                 │
├──────────────────────────────────────────────────────────────────┤
│ serverds          inetclient         demo_milestone1             │
│ (SOAP Server)     (SOAP Client)      (Demo App)                  │
└──────┬──────────────┬────────────────────┬──────────────────────┘
       │              │                    │
       └──────────────┼────────────────────┘
                      │
        ┌─────────────▼───────────────┐
        │   MODULES IMPLEMENTATE      │
        ├─────────────────────────────┤
        │ - server.c                  │
        │ - soapds.c (SOAP handlers)  │
        │ - processing.c              │
        │ - config_loader.c           │
        │ - client.c                  │
        └─────────────┬───────────────┘
                      │
        ┌─────────────▼───────────────┐
        │  GENERATED BY gSOAP         │
        ├─────────────────────────────┤
        │ - soapC.c                   │
        │ - soapServer.c              │
        │ - soapH.h                   │
        │ - proto.nsmap               │
        └─────────────┬───────────────┘
                      │
        ┌─────────────▼───────────────┐
        │  EXTERNAL LIBRARIES         │
        ├─────────────────────────────┤
        │ - gsoap                     │
        │ - libconfig                 │
        │ - librosacpp (LibrosaC)     │
        │ - libpthread                │
        │ - libm (math)               │
        │ - libc (POSIX)              │
        └─────────────────────────────┘
\end{verbatim}

% =============================================================================
\section{Header Files Structure}
% =============================================================================

\begin{lstlisting}
// include/server.h
#ifndef SERVER_H
#define SERVER_H
typedef struct { ... } server_config_t;
#define MAX_PATH_LEN 512
#define DEFAULT_PORT 8080
// ... constante
#endif

// include/processing.h
#ifndef PROCESSING_H
#define PROCESSING_H
typedef struct { ... } proc_request_t;
typedef struct { ... } proc_result_t;
int process_spectrogram(...);
#endif

// include/config_loader.h
#ifndef CONFIG_LOADER_H
#define CONFIG_LOADER_H
int  config_load(const char*path, server_config_t*out);
void config_free(server_config_t*cfg);
#endif

// src/proto.h
// gSOAP interface definition
//gSOAP ns serviciu
struct ns__AnalyzeAudioRequest { ... };
struct ns__AnalyzeAudioResponse { ... };
struct ns__JobStatusRequest { ... };
// ... mai multe structs
\end{lstlisting}

% =============================================================================
\section{Revision History}
% =============================================================================

\begin{table}[H]
    \centering
    \begin{tabular}{|l|c|l|}
        \toprule
        \textbf{Versiune} & \textbf{Data} & \textbf{Modificări} \\
        \midrule
        1.0 & 2026-04-21 & Versiune inițială completă --- SDD \\
        \bottomrule
    \end{tabular}
\end{table}

\end{document}
